% Generated draft: requires mathematical and source review.
\documentclass[11pt]{article}
\usepackage[margin=25mm]{geometry}
\usepackage{fontspec}
\setmainfont{Latin Modern Roman}
\usepackage{amsmath,amssymb,mathtools,unicode-math}
\setmathfont{Latin Modern Math}
\usepackage{enumitem,needspace,xurl,hyperref}
\setlength{\parindent}{0pt}
\setlength{\parskip}{5pt}
\setlength{\emergencystretch}{3em}
\newcommand{\N}{\mathbb{N}}
\newcommand{\Z}{\mathbb{Z}}
\newcommand{\Q}{\mathbb{Q}}
\newcommand{\R}{\mathbb{R}}
\newcommand{\C}{\mathbb{C}}
\newcommand{\F}{\mathbb{F}}
\newcommand{\A}{\mathbb{A}}
\newcommand{\PP}{\mathbb{P}}
\newcommand{\E}{\mathbb{E}}
\newcommand{\eps}{\varepsilon}
\newcommand{\dd}{\,\mathrm{d}}
\begin{document}
\section*{UnsolvedMath ID: 11000072}
\textbf{Source problem code:} catalogue entry
\subsection{Definitions and mathematical statement}
For every \(g\geq2\), the source asks:
\begin{enumerate}
\item Does \(\operatorname{Mod}_g\) have a finite-index subgroup \(H\) with an injective homomorphism \(H\to\operatorname{Homeo}(S^1)\)?
\item Does \(\operatorname{Mod}_g\) have a finite-index subgroup admitting an injective homomorphism into \(\operatorname{Homeo}(\R)\)?
\item Does \(\operatorname{Mod}_{g,1}\) have such a finite-index subgroup acting faithfully on \(\R\)?
\end{enumerate}
Faithful means injective, not free or properly discontinuous. The notation \(\operatorname{Mod}_{g,1}\) and its marked-point versus boundary convention are not defined in this extract. The printed targets allow all homeomorphisms; passing to the orientation kernel of a finite-index subgroup gives the corresponding existence question for orientation-preserving actions. For countable groups, faithful orientation-preserving actions on \(\R\) correspond to left-orderability, namely a total order invariant under left multiplication. Bi-orderability is a different property.

The source credits Thurston with braid-group orderability; reports that the Torelli groups \(I_g,I_{g,1}\), being residually torsion-free nilpotent, embed in \(\operatorname{Homeo}^+(\R)\); and says its extension (20) splits over \(I_{g,1}\). Residual torsion-free nilpotence means each nonidentity element survives in some torsion-free nilpotent quotient. It contrasts Witte [Wi], reporting no faithful circle or line action by any finite-index subgroup of \(\operatorname{Sp}(2g,\Z)\). The extension (20) is absent.

A following subsection supplies further context. Since \(\pi_1(\operatorname{Sp}(2g,\R))\cong\Z\), its universal covering Lie group gives a central extension; taking the inverse image of \(\operatorname{Sp}(2g,\Z)\) gives
\[
 1\longrightarrow\Z\longrightarrow
 \widetilde{\operatorname{Sp}}(2g,\Z)
 \longrightarrow\operatorname{Sp}(2g,\Z)\longrightarrow1.
\]
The source calls its defining class \(\zeta\in H^2(\operatorname{Sp}(2g,\Z);\Z)\) nontrivial and bounded and proportional to a K\"ahler class, and reports Deligne [De] proving the extension group not residually finite. It refers to a surjection from a preceding extension (19) and bounded cocycles for both, motivating a following question. Neither (19), the comparison map nor that question is supplied. The locally symmetric quotient's description as a \(K(\pi,1)\) also needs an orbifold or torsion-free-quotient convention; no such convention is silently inserted.
\subsection{Short English statement}
Ask whether finite-index subgroups of mapping class groups act faithfully on the circle or line, and retain the accompanying Torelli, symplectic-group and central-extension comparisons with their missing conventions.
\subsection{Sources}
\begin{itemize}
\item B. Farb, chapter in \emph{Problems on Mapping Class Groups and Related Topics}, Question 6.3, ``orderability,'' and following central-extension discussion, supplied PDF page 54. \url{https://www.math.uchicago.edu/~farb/papers/mcgbook.pdf}. Recorded access 2026-07-29; extraction pdf-font-structure.
\item Thurston, braid-orderability attribution; Witte [Wi], symplectic-group action obstruction; Deligne [De], non-residual finiteness. Complete bibliography entries are absent, as are extensions (19), (20) and the following question. The displayed symplectic extension is source (21).
\item Local UnsolvedMath registry, \texttt{data/unsolvedmath/problems.json}, permanent ID \(11000072\), exact source code \texttt{catalogue entry}; supplied OPEN-TRIAGE summary is tentative about the question and gives no precise subsequent reference.
\end{itemize}
\end{document}
